%% kapa: exam=1 round=30 year=2019 period=1 form=A subject=civil_law subject_ko=민법 seq=13 number=13
%% source: Q-Net 감정평가사 기출문제 (cst00309) artlSeq=5205205 file=extracted/1008/5205205/2019년 제30회 감정평가사 1차시험 문제지/2019년 제30회 감정평가사 1차시험 문제지/2019년도 제30회 감정평가사 1차 1교시 A형.hwp
%% conversion: kordoc 4.24.51 -> tools/gen1.py (auto); review=proofread-sonnet; flags=none
\begin{question}{13}
의사표시의 효력발생에 관한 설명으로 옳지 않은 것은? (다툼이 있으면 판례에 따름)
\begin{choices}
  \item 의사표시의 도달이란 상대방이 그 내용을 안 것을 의미한다.
  \item 의사표시의 부도달로 인한 불이익은 표의자가 부담한다.
  \item 도달주의의 원칙을 정하는 민법 제111조는 임의규정이므로 당사자는 약정으로 의사표시의 효력발생시기를 달리 정할 수 있다.
  \item 매매계약 승낙의 의사표시를 발신한 후 승낙자가 사망하였다고 하더라도 그 의사표시가 청약자에게 정상적으로 도달하였다면 매매계약은 유효하게 성립한다.
  \item 제한능력자는 원칙적으로 의사표시의 수령무능력자이다.
\end{choices}
\end{question}
\answer{1}
